\documentclass[10pt,a4paper]{amsart}
\usepackage[margin=19mm]{geometry}
\usepackage{amsmath,amssymb,mathtools}
\usepackage[scr=rsfs,cal=euler]{mathalfa}
\usepackage{xfrac}
\usepackage[unicode,hidelinks,bookmarks=false]{hyperref}
\newcommand{\ie}{i.e.\ }
\newcommand{\cf}{cf.\ }
\newcommand{\resp}{respectively\ }
\newcommand{\Subsection}{\subsection}
\providecommand{\nobreakdash}{\nobreak\hbox{-}}
\setlength{\emergencystretch}{2em}
\multlinegap0pt
\usepackage{bm}
\usepackage{bbm}
\usepackage{dsfont}
\usepackage{xspace}
\usepackage{cancel}
\usepackage{mathtools}
\usepackage{amsthm}
\makeatletter
\@ifundefined{lemma}{\newtheorem{lemma}{Lemma}}{}
\makeatother
\title[A Bi-Objective Markov Decision Process Design approach to re]{A Bi-Objective Markov Decision Process Design approach to redundancy allocation with dynamic maintenance for a parallel system}
\author{Luke Fairley, Rob Shone, Peter Jacko, Jefferson Huang}
\date{Source: arXiv:2504.12944v2 (CC BY 4.0)}
\begin{document}
\maketitle

\noindent\emph{Source and licence.} A Bi-Objective Markov Decision Process Design approach to redundancy allocation with dynamic maintenance for a parallel system. Version arXiv:2504.12944v2 (CC BY 4.0), distributed under the Creative Commons Attribution 4.0 International licence, as recorded in the arXiv licence metadata for this exact version and shown on the abstract page https://arxiv.org/abs/2504.12944v2. Full rights notice: licenses/arxiv-2504.12944-CC-BY-4.0.txt.\par\smallskip

\noindent\textit{Editorial scope.} Counted unit: the Lemma whose source label is \texttt{lemma:supermodular} in the article (source0.tex lines 1057--1060, source label \texttt{lemma:supermodular}), delivered together with the article's own complete original proof of it (source0.tex lines 1061--1109, one \texttt{proof} environment of 49 lines). No mathematics is written, repaired or re-derived: statement and proof are the article's own text line for line. Required context retained uncounted: none. Every definition, display and label of those lines is retained unchanged. Omitted from this package and not redistributed: 1--1056, 1110--1449. Nothing inside a retained excerpt is omitted or altered: every retained body is delivered exactly as the article's lines are written, so no omission receipt of any kind accompanies this package. Citations: the retained excerpts carry no citation command at all, so no bibliography entry of the article is transcribed and no reference is redistributed. Reference adaptation: none was needed and none was made; every reference inside the delivered unit resolves inside it, so no unresolved reference or citation is produced. Printed slips kept literal: no printed word, symbol, hypothesis or number of the counted statement or of its proof was altered. Layout and class: the article's own class is \texttt{article} with packages including amsmath, amsthm, amsfonts, graphicx, booktabs, color, natbib, bm, multirow, enumerate, subcaption, rotating, xcolor, tikz-cd; the wrapper is the lane's pinned \texttt{amsart} editorial wrapper at 10pt on A4, which restates the theorem-like environments the retained text needs and adds the article's own macro definitions that the retained text needs (none), with extra wrapper packages (bm, bbm, dsfont, xspace, cancel, mathtools). The delivered numbering is the unit's own. Assets: no retained span contains a figure environment, a graphics-inclusion command, a PDF-page inclusion, a table or a TikZ picture; no image file is copied into this package, the package declares no asset dependency and no image file is redistributed with it. Licence: version arXiv:2504.12944v2 of this article is distributed under the Creative Commons Attribution 4.0 International licence, as recorded in the arXiv licence metadata for this exact version and shown on the abstract page https://arxiv.org/abs/2504.12944v2. A full rights notice is delivered at licenses/arxiv-2504.12944-CC-BY-4.0.txt. Characters: every retained excerpt is pure ASCII, so the delivered engine, XeLaTeX with its default Latin Modern text font, reports no missing character.
\begin{lemma}
\label{lemma:supermodular}
    $g^{SET}$ is a supermodular set function.
\end{lemma}

\begin{proof}
We shall suppress the $SET$ superscript for brevity. Recall :
$$g(K) = \sum_{i\in\mathcal{L}'}q_ir_i + \sum_{i\in K}c_ip_i\prod_{j\in K,j\preceq i)}q_j + (1+\delta)c_N\prod_{i\in K}q_i.$$
Clearly, the first sum is a modular set function, so we need only show the supermodularity of the following:
$$f(K) := \sum_{i\in K}c_ip_i\prod_{j\in K,j\preceq i)}q_j + (1+\delta)c_N\prod_{i\in K}q_i,$$
as the positive linear combination of supermodular functions is supermodular.

To show supermodularity, for any subsets $X\subset Y \subset C'$, and any component $i\in \mathcal{L}'\backslash Y$, we must show that:
$$f(X\cup\{i\}) - f(X) \leq f(Y\cup\{i\}) - f(Y).$$

To simplify notation, we may use $K^{>i}$ and $K^{<i}$ to denote the partitions of a set $K$ with only the components greater than $l$ under the partial ordering, and the components considered lesser than $i$, respectively. We may also write:
$$f(i|K) := f(K\cup\{i\}) - f(K),$$
and $f(K;c)$ to be $f$ with $p$ replaced by $c$.

We first calculate $f(K\cup \{i\})$ as follows:
\begin{align*}
    f(K\cup\{i\}) &= \sum_{j\in K\cup\{i\}}\left( c_j p_j \prod_{k\in K\cup\{i\}, k<j}q_k\right) + (1+\delta)c_N\prod_{j\in K\cup\{i\}}q_j\\
    &= \sum_{j\in K, j < i}\left( c_j p_j \prod_{k\in K, k<j}q_k\right)\\
    &\quad+ c_ip_i\prod_{j < i}q_j
    + \sum_{j\in K, j > i}\left( c_j p_j \prod_{k\in K \cup \{i\}, k<j}q_k\right)\\
    &\quad+ (1+\delta)c_N \prod_{j\in K \cup \{i\}}q_j\\
    &= \sum_{j\in K, j < i}\left( c_j q_j \prod_{k\in K, k<j}q_k\right)\\
    &\quad+ c_ip_i\prod_{j < i}q_j
    + q_i\sum_{j\in K, j > i}\left( c_j p_j \prod_{k\in K, k<j}q_k\right)\\
    &\quad+ q_i(1+\delta)c_N \prod_{j\in K}q_j\\\vspace{1em}
    &= c_lx_1^l\prod_{k < l}(1 - x_1^k)\\
    &\quad+x_1^l\sum_{j\in K, j < l}\left( c_j x_1^j \prod_{k\in K, k<j}(1 - x_1^k)\right)\\
    &\quad+(1 - x_1^l)\left\{\sum_{j\in K}\left( c_j x_1^j \prod_{k\in K, k<j}(1 - x_1^i)\right) + p \prod_{i\in K}(1 - x_1^i)\right\}\\\vspace{1em}
    &= x_1^l\left\{c_l\prod_{k < l}(1 - x_1^k) + \sum_{i\in K, i < l}\left( c_i x_1^i \prod_{k\in K, k<j}(1 - p_k)\right)\right\}\\
    &\quad+ q_if(K)\\\vspace{1em}
    &= p_i f(K^{<i};c_i) + q_if(K).
\end{align*}
This makes intuitive sense. By including an extra component $i$, with probability $p_i$ you are at worst using component $i$ and ignoring all links worse than $i$, and with probability $q_i$ the link is unavailable and the system falls back onto the original links in $K$. Notably, this shows that $f$ is monotone. We then see that:
\begin{align*}
    f(i|K) &= f(K\cup\{i\}) - f(K)\\
    &= p_i\left(f(K^{<l};c_I) - f(K)\right).
\end{align*}
As the constant $p_i$ does not depend on $K$, we need only focus on $f(K^{<l};c_l) - f(K)$. We see:
\begin{align*}
    f(K^{<i};c_i) - f(K) &= c_i\prod_{k\in K, k < i}q_k + \sum_{j\in K, j < i}\left( c_j p_j \prod_{k\in K, k<j}q_k\right)\\
    &\quad-\left\{\sum_{j\in K}\left( c_j p_j \prod_{k\in K, k<j}q_j\right) + p \prod_{j\in K}q_j\right\}\\
    &= c_i\prod_{k\in K, k < i}q_k - \sum_{j\in K, j > i}\left( c_j p_j \prod_{k\in K, k<j}q_k\right) - p\prod_{j\in K}q_j\\
    &= \left\{\prod_{k\in K, k < i}q_k\right\}\left\{ c_i - \sum_{j\in K, j > i}c_j p_j \prod_{k\in K, i < k < j}q_k - p\prod_{j\in K, j > i}q_k\right\}\\
    &= \left\{\prod_{k\in K, k < i}q_k\right\}\left\{c_i - f(K^{>i})\right\}\\
    &=-\left\{\prod_{k\in K, k < i}q_k\right\}\left\{f(K^{>i}) - c_i\right\}.
\end{align*}

The magnitude of this expression decreases as we add more links to $K$ that are cheaper than $i$ by the first term, and decreases as we add more links more expensive than $i$ by the second term, as $f(K^{>i}) \geq c_i$. Therefore the function is supermodular. 
\end{proof}

\end{document}
